\documentclass[a4paper, twoside, 12pt]{report}

\usepackage[spanish, es-tabla]{babel}
\usepackage[utf8]{inputenc}
\usepackage[T1]{fontenc}
\usepackage{booktabs}
\usepackage{placeins}
\usepackage{graphicx}
\usepackage{float}
\usepackage{subcaption}
\usepackage{array}
\usepackage{multirow}
\usepackage{longtable}
\usepackage[a4paper,top=3cm,bottom=2cm,left=3cm,right=3cm,marginparwidth=1.75cm]{geometry}
\usepackage{amsmath,amssymb}
\usepackage{siunitx}
\usepackage{parskip}
\usepackage[colorinlistoftodos]{todonotes}
\usepackage[colorlinks=true, allcolors=blue]{hyperref}
\hypersetup{linkcolor=black,urlcolor=blue,citecolor=black}
\usepackage{fancyhdr}
\pagestyle{fancy}
\fancyhf{}
\fancyhead[L]{\nouppercase{\leftmark}}
\fancyhead[R]{\thepage}

% Reduce espacio vertical antes del titulo de capitulo
\makeatletter
\def\@makechapterhead#1{%
  \vspace*{20\p@}%
  {\parindent \z@ \raggedright \normalfont
    \ifnum \c@secnumdepth >\m@ne
      \huge\bfseries \@chapapp\space \thechapter
      \par\nobreak
      \vskip 20\p@
    \fi
    \interlinepenalty\@M
    \Huge \bfseries #1\par\nobreak
    \vskip 30\p@
  }}
\def\@makeschapterhead#1{%
  \vspace*{20\p@}%
  {\parindent \z@ \raggedright
    \normalfont
    \interlinepenalty\@M
    \Huge \bfseries  #1\par\nobreak
    \vskip 30\p@
  }}
\makeatother

\newcommand{\doccode}{INTI-REP-PLD-2025}
\newcommand{\doctitle}{Sistema de Microscopia y caracterizacion de muestras}
\newcommand{\docsubtitle}{Payload Microscopia}
\newcommand{\docproject}{Proyecto INTISAT -- APSCO}
\newcommand{\docauthor}{Jamin Yauri}
\newcommand{\docversion}{v1.0}
\newcommand{\docdate}{\today}
\newcommand{\HRule}{\rule{\linewidth}{0.5mm}}

\begin{document}

\begin{titlepage}
    \centering
    \vspace*{-1.5cm}
    \noindent
    \begin{minipage}{0.15\textwidth}
        
\includegraphics[width=2cm]{../../Images/Logo_APSCO.png}
    \end{minipage}%
    \hfill
    \begin{minipage}{0.8\textwidth}
        \raggedleft
        \Huge \textbf{APSCO}
    \end{minipage}

    \vspace{0.2cm}
    \HRule \\[0.5cm]

    \textbf{\LARGE \docproject}\\[0.5cm]
    \textsc{\Large Reporte Tecnico: \doccode}\\[0.3cm]
    \textsc{\Large Sistema: Payload Microscopia}\\[0.3cm]

    {\Large \docdate}\\[0.5cm]

    \HRule \\[0.4cm]
    { \huge \bfseries \doctitle: \\ \Large \docsubtitle}\\[0.4cm]
    \HRule \\[0.5cm]

    \noindent
    \begin{flushleft}
        \large \textbf{Preparado por:} \\
        \large Instituto de Radioastronomia (INRAS) - Equipo INTISAT\\[0.3cm]

        \large \textbf{Autor(es):} \large \docauthor\\[0.3cm]

        \large \textbf{Supervisado por:} \large M.Sc. Neils Edison Vílchez Lagos\\[0.3cm]

        \large \textbf{Version:} \docversion \\[0.1cm]
        \large \textbf{Estado:} Reporte Tecnico
    \end{flushleft}

    \vfill

    \begin{center}
        
\includegraphics[width=8cm]{../../Images/Intisat_logo.png}
    \end{center}

    \vspace{0.1cm}
\end{titlepage}

\pagenumbering{roman}

\begin{abstract}
    \addcontentsline{toc}{chapter}{Resumen Ejecutivo}
Este documento presenta los avances técnicos en el desarrollo, calibración y vali-
dación metrológica de un sistema de microscopía computacional basado en sensor
CMOS, concebido como una carga útil experimental (payload) para el proyecto IN-
TISAT, bajo la supervisión del Instituto de Radioastronomía (INRAS).

El payload tiene como objetivo demostrar la viabilidad de un sistema de caracteri-
zación de muestras biológicas y biomiméticas mediante una arquitectura compacta,
modular y de bajo consumo, adecuada para plataformas con recursos limitados. Pa-
ra ello, se emplea una estrategia de adquisición de imágenes basada en hardware
mínimo y procesamiento computacional, orientada a la extracción de información
morfológica a escala micrométrica.

Como etapa fundamental del desarrollo, el presente trabajo se enfoca en la validación
de la escala espacial (µm/pixel) del sistema de imagen. Se utilizan microesferas de
poliestireno de diámetro nominal de 2 µm como objetos patrón, permitiendo contras-
tar las mediciones obtenidas por el sistema CMOS con mediciones independientes
realizadas mediante microscopía óptica convencional y microscopía electrónica de
barrido.

La validación experimental busca asegurar la coherencia entre las dimensiones físi-
cas reales y las mediciones digitales, considerando las limitaciones impuestas por la
resolución del sensor, la óptica empleada y los métodos de adquisición. Este pro-
ceso constituye un paso crítico previo a la caracterización de muestras biológicas
reales, garantizando la confiabilidad del sistema como herramienta de observación
micrométrica en un contexto experimental espacial.

Finalmente, se documentan las pruebas realizadas, los resultados preliminares ob-
tenidos y las líneas de trabajo planificadas para las siguientes etapas, orientadas
a la integración progresiva del payload y a la ampliación de sus capacidades de
caracterización.\end{abstract}

\chapter*{Historial de Revisiones}
\begin{center}
\begin{tabular}{p{2cm} p{3cm} p{9cm}}
\toprule
\textbf{Version} & \textbf{Fecha} & \textbf{Descripcion} \\
\midrule
v1.0 & \today & Reporte tecnico inicial. \\
\bottomrule
\end{tabular}
\end{center}

\clearpage
\tableofcontents
\listoffigures
\listoftables
\clearpage

\pagenumbering{arabic}

\chapter{Introduccion}
\section{Contexto y motivacion}
El desarrollo de instrumentos científicos compactos, de bajo consumo y alta funcio-
nalidad constituye un desafío clave para las misiones espaciales modernas, especial-
mente en plataformas con recursos limitados como los nanosatélites tipo CubeSat.
En este contexto, la microscopía computacional surge como una alternativa viable
frente a los sistemas ópticos tradicionales, al trasladar parte de la complejidad del
hardware hacia el procesamiento digital de la información.

El proyecto INTISAT, desarrollado bajo el marco institucional del Instituto de Ra-
dioastronomía (INRAS), incorpora una carga útil experimental orientada a la valida-
ción de tecnologías emergentes aplicables a observación y caracterización de muestras
a pequeña escala. Dentro de este marco, el payload de microscopía computacional
busca explorar la factibilidad de realizar observaciones micrométricas funcionales
empleando sensores CMOS comerciales, óptica mínima y algoritmos de procesa-
miento computacional.

A diferencia de los microscopios de laboratorio convencionales, el sistema desarrolla-
do no persigue fines clínicos ni de alta precisión metrológica absoluta, sino la demos-
tración tecnológica de un sistema modular capaz de adquirir información morfológica
básica bajo restricciones severas de masa, volumen, potencia y complejidad.

Un aspecto crítico del desarrollo de sistemas de imagen no comerciales es la au-
sencia de calibración metrológica de fábrica. En consecuencia, resulta indispensable
validar que las mediciones digitales obtenidas a partir de las imágenes capturadas
correspondan de manera coherente con dimensiones físicas reales. Esta validación
constituye un requisito previo para cualquier intento de caracterización de muestras
biológicas o biomiméticas mediante el payload propuesto.

En respuesta a esta necesidad, el presente trabajo se enfoca en la validación experi-
mental de la escala espacial del sistema de microscopía, utilizando objetos patrón y
técnicas de comparación cruzada con sistemas de microscopía calibrados.
\section{Objetivo general}
Demostrar la viabilidad de un sistema de microscopía computacional basado en
sensor CMOS como carga útil experimental para el proyecto INTISAT, mediante la
validación metrológica de su escala espacial y la verificación de la coherencia entre
las mediciones digitales obtenidas y las dimensiones físicas reales de objetos patrón
micrométricos.

\section{Objetivos Específicos}

Los objetivos específicos del presente trabajo son:

\begin{itemize}
    \item Diseñar e implementar un sistema de adquisición de imágenes basado en sensor CMOS y configuración óptica mínima, adecuado para su integración como payload experimental dentro del proyecto INTISAT.
    
    \item Establecer una metodología de validación metrológica de la escala espacial del sistema de imagen (µm/pixel), empleando objetos patrón micrométricos de dimensiones conocidas.
    
    \item Evaluar el desempeño del sistema frente a objetos patrón micrométricos, específicamente microesferas de poliestireno de diámetro nominal de 2 µm.
    
    \item Comparar las mediciones obtenidas mediante el sistema CMOS con mediciones independientes realizadas por microscopía óptica calibrada y microscopía electrónica de barrido (SEM), a fin de determinar la coherencia dimensional del sistema.
    
    \item Analizar las limitaciones físicas y computacionales del sistema, considerando efectos de resolución espacial, difracción, tamaño de pixel y procesamiento digital de imagen.
    
    \item Documentar los resultados experimentales obtenidos y establecer lineamientos técnicos para etapas futuras de desarrollo, optimización y ampliación de capacidades del payload.
\end{itemize}
\section{Alcance del Trabajo}

El presente trabajo comprende la validación experimental y documental de la escala espacial del sistema de microscopía computacional desarrollado como carga útil experimental del proyecto INTISAT, estableciendo las bases técnicas necesarias para su uso en futuras aplicaciones de caracterización de muestras biológicas y biomiméticas.

En particular, el estudio abarca:

\begin{itemize}
    \item La caracterización dimensional de objetos patrón micrométricos mediante el sistema CMOS desarrollado.
    
    \item La comparación de las mediciones obtenidas con técnicas de microscopía óptica convencional y microscopía electrónica de barrido, como referencia metrológica independiente.
    
    \item El análisis de las limitaciones físicas y computacionales del sistema en términos de resolución efectiva, efectos de difracción, tamaño de pixel y capacidad de observación morfológica.
    
    \item La evaluación preliminar del potencial del sistema para la observación de estructuras celulares de tamaño micrométrico, sin fines diagnósticos ni clínicos.
\end{itemize}

No forman parte del alcance del presente trabajo:

\begin{itemize}
    \item La validación clínica, diagnóstica o terapéutica de muestras biológicas.
    
    \item El uso del sistema como herramienta médica certificada.
    
    \item La interpretación biomédica de resultados asociados a patologías específicas.
    
    \item La implementación completa de técnicas avanzadas de reconstrucción computacional, tales como la microscopía ptychográfica de Fourier, como herramienta primaria de calibración metrológica.
\end{itemize}

Sin embargo, se reconoce que, una vez validada la coherencia metrológica del sistema y optimizadas sus capacidades de adquisición y procesamiento, el payload podría emplearse como plataforma experimental para la observación de líneas celulares ampliamente utilizadas en investigación (por ejemplo, células HeLa), exclusivamente con fines exploratorios, educativos o de investigación básica.

El presente trabajo se concibe como una etapa preliminar dentro del desarrollo progresivo del payload, proporcionando una base técnica sólida para futuras fases orientadas a la ampliación de capacidades de caracterización biológica y computacional.

\chapter{Marco Teórico}

\section{Microscopía y caracterización a escala micrométrica}

La microscopía constituye un conjunto de técnicas orientadas a la observación y caracterización de estructuras cuyo tamaño se encuentra por debajo del límite de resolución del ojo humano. En el contexto de la ingeniería y la investigación científica, la microscopía no se limita a la visualización cualitativa, sino que se emplea como herramienta para la medición, comparación y análisis morfológico de estructuras a escala micrométrica y submicrométrica.

En aplicaciones biológicas y biomiméticas, la microscopía permite identificar características estructurales tales como tamaños celulares, bordes, distribuciones espaciales y relaciones geométricas básicas. Estas observaciones constituyen una base fundamental para estudios exploratorios, educativos y de investigación básica, sin implicar necesariamente un uso clínico o diagnóstico.

Desde un punto de vista instrumental, los sistemas de microscopía pueden clasificarse según su arquitectura óptica, el tipo de detector empleado y el grado de procesamiento computacional involucrado en la formación de la imagen. En este marco, los sistemas tradicionales basados en lentes conviven con enfoques alternativos que priorizan la simplicidad del hardware y delegan parte de la interpretación visual al procesamiento posterior.

\section{Fundamentos físicos de la microscopía sin lentes}

La microscopía sin lentes (\textit{lensless microscopy}) es un enfoque en el cual la formación de imagen no depende de sistemas ópticos refractivos tradicionales, sino de la interacción directa entre la luz y la muestra, registrada por un detector fotosensible ubicado a corta distancia.

En configuraciones de microscopía de contacto, la muestra se coloca a una distancia del orden de micrómetros del plano del sensor, de modo que la imagen resultante está dominada por fenómenos de absorción, dispersión cercana y difracción. En este régimen, la luz transmitida o dispersada por la muestra proyecta patrones de intensidad directamente sobre el sensor, los cuales contienen información estructural de la muestra.

La ausencia de lentes reduce significativamente la complejidad mecánica y óptica del sistema, permitiendo arquitecturas compactas, robustas y de bajo consumo energético. Sin embargo, esta simplificación introduce limitaciones inherentes relacionadas con la resolución efectiva, la sensibilidad al acoplamiento óptico muestra–sensor y la influencia de efectos difractivos.

\section{Sensor CMOS como sistema de adquisición}

Los sensores CMOS (Complementary Metal-Oxide-Semiconductor) se han consolidado como una tecnología madura y ampliamente disponible para la adquisición de imágenes digitales. Su bajo consumo de potencia, alta integración y compatibilidad con sistemas embebidos los convierten en una opción adecuada para plataformas con recursos limitados.

En el contexto de microscopía sin lentes, el sensor CMOS actúa como el plano de detección donde se registra la distribución espacial de intensidad luminosa generada por la interacción de la luz con la muestra. Cada píxel del sensor integra fotones incidentes durante un intervalo de tiempo definido, produciendo una señal digital proporcional a la intensidad local.

El tamaño físico del píxel, la eficiencia cuántica, el ruido electrónico y la resolución nativa del sensor establecen límites fundamentales a la capacidad del sistema para resolver estructuras pequeñas.

\begin{table}[H]
\centering
\caption{Características técnicas relevantes del sensor CMOS utilizado}
\begin{tabular}{|l|l|}
\hline
\textbf{Parámetro} & \textbf{Valor típico} \\
\hline
Tipo de sensor & CMOS (Commercial Off-The-Shelf) \\
Resolución nativa & 2592 × 1944 píxeles \\
Tamaño de píxel & $\sim 1.4\,\mu m$ \\
Área activa & $\sim 3.7\,mm \times 2.7\,mm$ \\
Profundidad de bits & 10–12 bits \\
Control de exposición & Manual \\
Control de ganancia & Manual \\
Filtro de color & Matriz Bayer (RGB) \\
Interfaz de datos & CSI / interfaz digital directa \\
\hline
\end{tabular}
\end{table}

\section{Resolución, difracción y límites físicos del sistema}

La resolución espacial de un sistema de imagen se define como la capacidad de distinguir dos estructuras cercanas como entidades separadas. En microscopía sin lentes, la resolución efectiva está influenciada por:

\begin{itemize}
    \item Tamaño físico del píxel del sensor.
    \item Distancia muestra–sensor.
    \item Longitud de onda de la iluminación.
    \item Respuesta espacial del detector.
\end{itemize}

Cuando las dimensiones de la muestra son comparables al límite de resolución del sistema, las estructuras no se observan como objetos bien definidos, sino como patrones difusos dominados por fenómenos de difracción. Este comportamiento es esperable y responde a principios físicos fundamentales de propagación de la luz.

\section{Escala espacial y calibración dimensional}

La obtención de mediciones dimensionales a partir de imágenes digitales requiere establecer una relación confiable entre el espacio físico y el espacio discreto representado por los píxeles del sensor. Esta relación, comúnmente expresada como escala espacial en unidades de micrómetros por píxel ($\mu m/pixel$), debe validarse experimentalmente.

En desarrollos experimentales basados en sensores CMOS, la escala inicial puede estimarse a partir de parámetros nominales del sensor. No obstante, esta estimación debe contrastarse mediante procedimientos independientes antes de utilizar el sistema para caracterización dimensional de muestras reales.

La calibración dimensional constituye, por tanto, un paso crítico para asegurar la trazabilidad y coherencia metrológica del sistema.

\section{Objetos patrón micrométricos para validación}

Los objetos patrón micrométricos se emplean como referencias físicas conocidas para validar la escala espacial y la capacidad de medición de un sistema de imagen. Entre los más utilizados se encuentran microesferas poliméricas y patrones litográficos con geometrías definidas.

Las microesferas de poliestireno presentan diámetros nominales certificados y una dispersión de tamaños conocida, lo que permite evaluar la coherencia dimensional del sistema.

\begin{table}[H]
\centering
\caption{Características típicas de microesferas de poliestireno}
\begin{tabular}{|l|l|}
\hline
\textbf{Parámetro} & \textbf{Valor típico} \\
\hline
Material & Poliestireno \\
Diámetro nominal & 2 $\mu m$ \\
Coeficiente de variación & $\sim 5\%$ \\
Forma & Esférica \\
Medio de suspensión & Solución acuosa \\
Índice de refracción & $\sim 1.59$ \\
\hline
\end{tabular}
\end{table}

\begin{table}[H]
\centering
\caption{Comparación cualitativa de objetos patrón}
\begin{tabular}{|l|c|c|}
\hline
\textbf{Característica} & \textbf{Microesferas} & \textbf{Patrones litográficos} \\
\hline
Geometría & Esférica & Lineal / reticular \\
Dimensión conocida & Nominal (distribuida) & Certificada \\
Influencia de difracción & Alta cerca del límite & Menor \\
Robustez metrológica & Media & Alta \\
Adecuación para $\mu m/pixel$ & Complementaria & Principal \\
\hline
\end{tabular}
\end{table}

\section{Microscopía computacional como apoyo a la interpretación}

La microscopía computacional engloba técnicas orientadas a mejorar la interpretabilidad de imágenes mediante procesamiento posterior. Estas técnicas no sustituyen la adquisición física de la imagen, sino que actúan como herramientas complementarias para realce visual y análisis exploratorio.

En el presente trabajo, el procesamiento computacional se emplea como apoyo a la interpretación de estructuras ya presentes en la imagen base, reconociendo explícitamente que no incrementa la resolución óptica real del sistema ni genera información física adicional.

Este enfoque garantiza la trazabilidad dimensional y mantiene una separación clara entre adquisición física y postprocesamiento digital.

\chapter{Descripción del Sistema}

\section{Arquitectura General}

El sistema de microscopía desarrollado se concibe como una carga útil experimental compacta, modular y de bajo consumo, diseñada para su integración en el proyecto INTISAT. La arquitectura general prioriza la simplicidad funcional, la trazabilidad de los datos y la reproducibilidad experimental, en concordancia con las restricciones propias de plataformas tipo CubeSat.

Desde un punto de vista funcional, el sistema se organiza en cuatro bloques principales: adquisición óptica, iluminación, control y procesamiento, y almacenamiento de datos. Esta descomposición permite aislar las funciones físicas de captura de imagen de las funciones lógicas de control y análisis, facilitando tanto la verificación del sistema como su futura ampliación.

El sistema opera bajo una arquitectura centralizada, en la cual una computadora de carga útil coordina la adquisición de imágenes, el control de la iluminación y el registro de metadatos, asegurando coherencia entre las condiciones experimentales y los datos generados.

\begin{figure}[H]
    \centering
    \includegraphics[width=0.65\textwidth]{fig/Fig_03_01_Arquitectura_General_Payload.png}
    \caption{Arquitectura general del payload de microscopía: bloques funcionales y flujo de datos.}
    \label{fig:arch_general}
\end{figure}

\section{Bloque de adquisición óptica}

El bloque de adquisición óptica está compuesto por un sensor CMOS comercial operando en configuración de microscopía sin lentes. En esta arquitectura, la muestra se coloca en contacto cercano con el plano del sensor, eliminando el uso de elementos ópticos refractivos para la formación de imagen.

La imagen se genera a partir de la interacción directa entre la luz incidente y la muestra, siendo registrada como una distribución espacial de intensidad en la matriz de píxeles del sensor. El bloque permite el control manual de parámetros críticos como el tiempo de exposición y la ganancia, deshabilitando mecanismos automáticos para garantizar consistencia entre capturas consecutivas.

Este enfoque permite una arquitectura óptica simple, robusta y adecuada para validaciones experimentales, aunque introduce limitaciones asociadas a la resolución efectiva y a la sensibilidad al acoplamiento óptico muestra--sensor.

\begin{figure}[H]
    \centering
    \includegraphics[width=0.92\textwidth]{fig/Fig_03_02_Lensless_Setup_CMOS_Muestra.png}
    \caption{Esquema conceptual de microscopía sin lentes: acoplamiento muestra--sensor y registro directo del patrón de intensidad.}
    \label{fig:lensless_setup}
\end{figure}

\section{Bloque de iluminación}

El bloque de iluminación proporciona la fuente de luz necesaria para la observación de la muestra y está implementado mediante un sistema plano y programable. Este bloque permite generar patrones de iluminación controlados en intensidad y geometría, adaptándose tanto a la operación nominal como a modos experimentales.

En condiciones nominales, el sistema emplea iluminación aproximadamente uniforme con el objetivo de maximizar el contraste de la muestra y reducir artefactos asociados a variaciones espaciales de intensidad. De manera complementaria, se dispone de la capacidad de generar patrones estructurados de iluminación para evaluación experimental de técnicas computacionales avanzadas.

El diseño del bloque de iluminación prioriza la estabilidad temporal, la reproducibilidad de los patrones generados y la integración compacta con el resto del sistema.

\begin{figure}[H]
    \centering
    \includegraphics[width=0.95\textwidth]{fig/Fig_03_03_Patrones_Iluminacion_Uniforme_Rejilla_Puntual.png}
    \caption{Modos de iluminación evaluados: (a) uniforme, (b) estructurada tipo rejilla y (c) puntual/dirigida.}
    \label{fig:illum_patterns}
\end{figure}

\section{Bloque de control y procesamiento}

El bloque de control y procesamiento constituye el núcleo lógico del sistema y se encarga de coordinar la operación de todos los subsistemas. Este bloque se implementa mediante una computadora de carga útil que ejecuta el software de control del payload.

Sus funciones principales incluyen la configuración del sensor CMOS, el control del sistema de iluminación, la gestión de los modos operativos y la ejecución de rutinas básicas de calibración, tales como la adquisición de imágenes de referencia para corrección de ruido y no uniformidades.

Asimismo, este bloque gestiona el registro automático de parámetros de adquisición y métricas básicas de calidad de imagen, garantizando la trazabilidad experimental de cada captura. El procesamiento computacional avanzado se concibe como una etapa complementaria y no afecta la operación nominal del sistema.

\subsection{Computadora de carga útil: Raspberry Pi 5}

La computadora de carga útil del sistema de microscopía está implementada mediante una Raspberry Pi 5, seleccionada como plataforma central de control, adquisición y gestión de datos del payload.

La Raspberry Pi 5 actúa como elemento coordinador de todos los subsistemas, encargándose de la configuración y control del sensor CMOS, la gestión del sistema de iluminación, la ejecución de los modos operativos y el almacenamiento de los datos generados. Su capacidad de procesamiento y su compatibilidad con interfaces de alta velocidad permiten una integración directa y estable con el sistema de adquisición óptica.

La elección de la Raspberry Pi 5 responde a criterios de disponibilidad, flexibilidad, bajo consumo relativo y amplia adopción en entornos de instrumentación experimental. Asimismo, su ecosistema de software y soporte para sistemas operativos basados en Linux facilita la implementación de rutinas de control, registro de metadatos y procesamiento básico de imágenes.

Desde el punto de vista arquitectónico, la Raspberry Pi 5 se considera un componente funcional del payload y no una plataforma de propósito general. Su operación se limita estrictamente a las tareas necesarias para la adquisición, control y trazabilidad experimental del sistema de microscopía.

\begin{table}[H]
\centering
\caption{Características técnicas relevantes de la Raspberry Pi 5 como computadora de carga útil}
\begin{tabular}{|l|l|}
\hline
\textbf{Parámetro} & \textbf{Valor relevante} \\
\hline
Modelo & Raspberry Pi 5 \\
Arquitectura & ARM de 64 bits \\
Rol en el sistema & Control y adquisición \\
Interfaces utilizadas & CSI (cámara), I2C (iluminación), GPIO \\
Sistema operativo & Linux embebido \\
Capacidad de procesamiento & Adecuada para control y análisis básico \\
Almacenamiento & Memoria externa (microSD / SSD) \\
\hline
\end{tabular}
\end{table}

\section{Bloque de almacenamiento y datos}

El bloque de almacenamiento se encarga de la persistencia de la información generada durante la operación del sistema. Incluye tanto las imágenes adquiridas como los metadatos asociados a cada captura.

Las imágenes se almacenan en formatos digitales sin compresión o con compresión mínima, preservando la información original capturada por el sensor. De manera complementaria, se generan archivos de metadatos estructurados que contienen parámetros de adquisición, configuraciones de iluminación y marcas temporales.

Esta organización de datos permite una trazabilidad completa del proceso experimental y facilita el análisis posterior en etapas de validación y procesamiento.

\begin{figure}[H]
    \centering
    \includegraphics[width=0.92\textwidth]{fig/Fig_03_04_Estructura_Datos_Imagen_Metadatos.png}
    \caption{Estructura de datos propuesta: imágenes + metadatos para trazabilidad experimental.}
    \label{fig:data_schema}
\end{figure}

\section{Flujo funcional del sistema}

El flujo funcional del sistema durante una operación nominal de captura puede resumirse de la siguiente manera:

\begin{enumerate}
    \item Inicialización del sistema de control y configuración del sensor CMOS.
    \item Configuración del patrón e intensidad del sistema de iluminación.
    \item Verificación de parámetros de adquisición y condiciones de exposición.
    \item Captura de la imagen por parte del sensor CMOS.
    \item Registro de métricas básicas de calidad de imagen.
    \item Almacenamiento de la imagen y generación de metadatos asociados.
\end{enumerate}

Este flujo se mantiene constante para la operación nominal del payload, variando únicamente en modos experimentales el número de capturas y la secuencia de iluminación aplicada. La separación clara entre adquisición y procesamiento garantiza que la imagen base represente fielmente la información física capturada por el sistema.

\begin{figure}[H]
    \centering
    \includegraphics[width=0.95\textwidth]{fig/Fig_03_05_Flujo_Funcional_Captura.png}
    \caption{Flujo funcional nominal: control, iluminación, captura, métricas y almacenamiento.}
    \label{fig:functional_flow}
\end{figure}

\begin{table}[H]
\centering
\caption{Resumen de bloques funcionales del sistema}
\begin{tabular}{|l|l|}
\hline
\textbf{Bloque} & \textbf{Función principal} \\
\hline
Adquisición óptica & Captura directa mediante sensor CMOS \\
Iluminación & Iluminación uniforme y estructurada \\
Control y procesamiento & Coordinación y registro de parámetros \\
Almacenamiento y datos & Persistencia de imágenes y metadatos \\
\hline
\end{tabular}
\end{table}

\chapter{Metodología}

\section{Enfoque general de validación}

La metodología del presente trabajo se fundamenta en una validación metrológica preliminar orientada a verificar la coherencia entre mediciones digitales obtenidas a partir del sistema de microscopía basado en sensor CMOS y dimensiones físicas reales de objetos patrón micrométricos.

Dado que el sistema de imagen es un desarrollo experimental y no cuenta con calibración de fábrica, la validación se realiza mediante una estrategia de comparación cruzada. Esta estrategia consiste en medir los mismos objetos patrón utilizando: (i) el sistema CMOS propio, (ii) un sistema de microscopía óptica calibrada, y (iii) microscopía electrónica de barrido (SEM), con el fin de contrastar resultados y establecer una referencia dimensional independiente.

\begin{figure}[H]
    \centering
    \includegraphics[width=0.65\textwidth]{fig/Fig_04_01_Esquema_Validacion_Cruzada.png}
    \caption{Esquema de validación cruzada: comparación entre el sistema CMOS propio, microscopía óptica calibrada y SEM.}
    \label{fig:cross_validation}
\end{figure}

\section{Materiales y equipos}

Para la ejecución del procedimiento experimental se emplearon los materiales y equipos listados a continuación:

\begin{itemize}
    \item Sistema de microscopía basado en sensor CMOS en configuración \textit{lensless/contact}.
    \item Computadora de carga útil: Raspberry Pi 5 (control de sensor, iluminación y almacenamiento).
    \item Sistema de iluminación programable del payload (modo nominal y experimental).
    \item Objetos patrón micrométricos: microesferas de poliestireno de diámetro nominal $2\,\mu m$.
    \item Acceso a microscopía óptica calibrada en laboratorio (PUCP).
    \item Acceso a microscopía electrónica de barrido (SEM) en el Centro de Caracterización de Materiales (CAM PUCP).
\end{itemize}

\begin{table}[H]
\centering
\caption{Resumen de equipos y su rol en la validación}
\begin{tabular}{|l|l|}
\hline
\textbf{Elemento} & \textbf{Rol en la metodología} \\
\hline
Sistema CMOS propio & Adquisición principal y medición digital \\
Raspberry Pi 5 & Control de captura, metadatos y almacenamiento \\
Iluminación programable & Control de condiciones de iluminación \\
Microscopía óptica calibrada & Referencia dimensional independiente \\
SEM (CAM PUCP) & Confirmación dimensional/morfológica de alta resolución \\
\hline
\end{tabular}
\end{table}

\section{Preparación de objetos patrón}

Las microesferas de poliestireno de $2\,\mu m$ se emplean como objetos patrón preliminares debido a su disponibilidad y uso extendido como referencia micrométrica. Las microesferas se encuentran en suspensión acuosa, por lo que se requiere preparar una muestra adecuada para adquisición en configuración de microscopía de contacto.

De manera general, la preparación de la muestra considera:

\begin{itemize}
    \item Homogeneización suave de la suspensión para reducir efectos de sedimentación.
    \item Depósito controlado de un volumen reducido de suspensión sobre un soporte compatible.
    \item Implementación de una interfaz muestra--sensor que permita contacto cercano sin contaminar el sensor.
    \item Registro de un identificador de muestra (\textit{Sample Tag}) para trazabilidad experimental.
\end{itemize}

\begin{figure}[H]
    \centering
    \includegraphics[width=0.92\textwidth]{fig/Fig_04_02_Preparacion_Muestra_Microesferas.png}
    \caption{Preparación conceptual de la muestra: suspensión, depósito controlado y configuración muestra--sensor en modo contacto.}
    \label{fig:sample_prep}
\end{figure}

\section{Procedimiento de adquisición con el sistema CMOS}

La adquisición de imágenes se realiza mediante el sistema CMOS operando en modo nominal (\textit{Capture-Single}). El objetivo es obtener imágenes reproducibles, evitando cambios automáticos de parámetros del sensor.

La secuencia general de captura es la siguiente:

\begin{enumerate}
    \item Inicialización del sensor CMOS y del sistema de iluminación.
    \item Desactivación de mecanismos automáticos de exposición y balance (cuando aplique).
    \item Configuración manual de parámetros de adquisición (tiempo de exposición y ganancia).
    \item Selección del patrón de iluminación nominal (iluminación uniforme).
    \item Captura de imagen a resolución nativa del sensor.
    \item Cálculo de métricas básicas de calidad (histograma, saturación).
    \item Almacenamiento de la imagen y generación del archivo de metadatos asociado.
\end{enumerate}

Con el fin de asegurar consistencia, se realizan múltiples capturas consecutivas bajo la misma configuración. Para cada captura, el sistema registra metadatos experimentales que permiten reconstruir las condiciones exactas de adquisición.

\begin{table}[H]
\centering
\caption{Parámetros mínimos registrados por captura para trazabilidad}
\begin{tabular}{|l|l|}
\hline
\textbf{Parámetro} & \textbf{Descripción} \\
\hline
ExposureTime & Tiempo de integración del sensor \\
AnalogueGain & Ganancia analógica del sensor \\
Resolution & Resolución de captura \\
OLED Pattern ID & Identificador del patrón de iluminación \\
OLED Intensity & Intensidad relativa del patrón \\
Timestamp & Marca temporal de la captura \\
Saturation Ratio & Porcentaje de píxeles saturados \\
Sample Tag & Identificador experimental de la muestra \\
\hline
\end{tabular}
\end{table}

\section{Medición digital y estimación de escala espacial}

La medición digital consiste en estimar el tamaño aparente de los objetos patrón sobre la imagen capturada. Para ello se seleccionan regiones de interés (ROI) que contienen microesferas aisladas y se realiza la medición de su diámetro aparente en píxeles.

La conversión de píxeles a unidades físicas se realiza mediante una escala espacial ($\mu m/pixel$) definida inicialmente por estimación geométrica o por referencia preliminar. Posteriormente, esta escala se ajusta o valida mediante comparación con mediciones externas.

El procedimiento general de medición digital incluye:

\begin{itemize}
    \item Selección de ROI con microesferas distinguibles.
    \item Segmentación o detección del contorno aparente (según calidad de imagen).
    \item Medición del diámetro aparente en píxeles.
    \item Conversión a micrómetros mediante la escala ($\mu m/pixel$).
    \item Cálculo de estadísticos básicos (promedio y dispersión).
\end{itemize}

\begin{figure}[H]
    \centering
    \includegraphics[width=0.92\textwidth]{fig/Fig_04_03_Medicion_Digital_ROI_Diametro.png}
    \caption{Flujo de medición digital: selección de ROI, medición de diámetro aparente en píxeles y conversión a micrómetros.}
    \label{fig:digital_measurement}
\end{figure}

\section{Comparación cruzada con microscopía óptica y SEM}

Las microesferas medidas mediante el sistema CMOS se contrastan con mediciones obtenidas en microscopía óptica calibrada y microscopía electrónica de barrido (SEM). Esta comparación permite validar de manera independiente el diámetro real, la distribución de tamaños y la morfología general.

La comparación cruzada se realiza mediante:

\begin{itemize}
    \item Obtención de micrografías representativas en cada técnica.
    \item Medición del diámetro de un conjunto de microesferas en cada sistema.
    \item Comparación de promedios y dispersión entre técnicas.
    \item Identificación de discrepancias atribuibles a límites de resolución y efectos difractivos en el sistema CMOS.
\end{itemize}

\section{Criterios de calidad y control de errores}

Para asegurar confiabilidad en la validación, se emplean criterios mínimos de calidad durante la adquisición y el análisis:

\begin{itemize}
    \item \textbf{Control de saturación:} la imagen no debe presentar saturación significativa bajo condiciones nominales.
    \item \textbf{Consistencia entre capturas:} variación relativa baja de métricas globales (promedio e histograma).
    \item \textbf{Trazabilidad completa:} cada imagen debe contar con metadatos completos y un identificador de muestra.
    \item \textbf{Interpretación conservadora:} las mediciones se reportan reconociendo explícitamente limitaciones de resolución y difracción.
\end{itemize}

En caso de observar patrones dominados por difracción o falta de definición de bordes, la medición dimensional se considera preliminar y se prioriza la comparación con patrones litográficos o mediciones SEM como referencia de mayor robustez.

\section{Procedimiento de adquisición mediante control programático}

La adquisición de imágenes se realiza mediante control programático del sensor CMOS OV5647, utilizando una Raspberry Pi 5 como plataforma de ejecución. Este enfoque permite un control preciso y reproducible de los parámetros de captura, condición necesaria para la validación metrológica y la comparación entre múltiples adquisiciones.

El control programático garantiza que todas las imágenes se obtengan bajo configuraciones conocidas y constantes, evitando variaciones introducidas por algoritmos automáticos del sensor.

\subsection{Plataforma de control y entorno de ejecución}

La Raspberry Pi 5 actúa como computadora de carga útil del sistema, ejecutando un entorno Linux embebido y empleando una interfaz de control de cámara basada en la librería \textit{Picamera2}. Esta interfaz permite acceder de manera directa a parámetros internos del sensor OV5647 sin requerir manipulación de registros a bajo nivel.

El software de control inicializa el sensor, configura los modos de operación, ejecuta la secuencia de captura y gestiona el almacenamiento de imágenes y metadatos asociados. Todas las operaciones de adquisición se realizan en modo manual, deshabilitando mecanismos automáticos del sensor.

\subsection{Configuración del sensor CMOS}

Previo a cada sesión de adquisición, el sensor CMOS se configura mediante parámetros definidos explícitamente por software. Los parámetros críticos controlados durante la captura incluyen:

\begin{itemize}
    \item Tiempo de exposición (\textit{Exposure Time}), que determina el intervalo de integración en cada píxel.
    \item Ganancia analógica (\textit{Analog Gain}), aplicada antes de la conversión analógica-digital.
    \item Resolución de captura, definida por el tamaño activo del arreglo de píxeles.
    \item Desactivación de control automático de exposición (AE) y balance de blancos (AWB), para asegurar consistencia entre capturas.
\end{itemize}

La selección de estos parámetros se realiza considerando el compromiso entre nivel de señal, ruido electrónico y riesgo de saturación, priorizando siempre el incremento del tiempo de exposición y la iluminación antes de incrementar la ganancia analógica.

\subsection{Secuencia programática de captura}

La captura de una imagen sigue una secuencia programática definida, diseñada para asegurar estabilidad del sensor y reproducibilidad experimental:

\begin{enumerate}
    \item Inicialización del sensor CMOS y aplicación de la configuración de captura.
    \item Aplicación explícita de parámetros manuales de exposición y ganancia.
    \item Estabilización temporal del sensor para asegurar integración completa.
    \item Captura de la imagen a resolución nativa del sensor.
    \item Conversión de la imagen adquirida a escala de grises para análisis posterior (cuando aplique).
    \item Retorno a modo de espera para la siguiente captura.
\end{enumerate}

\subsection{Registro automático de metadatos}

Durante la adquisición, el sistema genera automáticamente metadatos asociados a cada imagen capturada. Estos incluyen información del sensor, parámetros de captura, configuración de iluminación y métricas básicas de calidad de imagen.

El registro automático de metadatos asegura la trazabilidad experimental de cada imagen, facilita la reproducibilidad del experimento y permite la comparación posterior con mediciones obtenidas mediante otras técnicas.

Cada imagen y su archivo de metadatos asociado se almacenan de manera conjunta, constituyendo una unidad experimental indivisible para las etapas posteriores de análisis y validación.

\begin{figure}[H]
    \centering
    \includegraphics[width=0.92\textwidth]{fig/Fig_04_04_Pipeline_Captura_Metadatos.png}
    \caption{Pipeline de captura programática: configuración manual, adquisición, cálculo de métricas y almacenamiento con metadatos.}
    \label{fig:pipeline_capture}
\end{figure}

\chapter{Resultados Experimentales}

\section{Validación dimensional independiente de objetos patrón}

Como paso previo a la evaluación del sistema CMOS, se realizó la caracterización dimensional independiente de las microesferas de poliestireno de diámetro nominal $2\,\mu m$, empleando microscopía óptica calibrada y microscopía electrónica de barrido (SEM).

Las imágenes obtenidas mediante SEM permitieron observar con alta definición la morfología esférica de las partículas y realizar mediciones directas del diámetro aparente.

\begin{figure}[H]
    \centering
    \includegraphics[width=0.5\textwidth]{fig/Fig_05_01_SEM_Microesferas.png}
    \caption{Micrografía SEM representativa de microesferas de poliestireno (2 µm nominal).}
    \label{fig:sem_microesferas}
\end{figure}

A partir de la medición de múltiples partículas, se obtuvo un diámetro promedio consistente con el valor nominal dentro del margen de dispersión esperado.

\begin{table}[H]
\centering
\caption{Resultados de medición dimensional mediante SEM}
\begin{tabular}{|c|c|}
\hline
\textbf{Número de muestras} & $n$ \\
Diámetro promedio & $\bar{d}_{SEM}$ \\
Desviación estándar & $\sigma_{SEM}$ \\
\hline
\end{tabular}
\end{table}

Estos resultados constituyen la referencia dimensional independiente utilizada para contrastar las mediciones obtenidas mediante el sistema CMOS.
\section{Adquisición directa mediante sensor CMOS en configuración lensless}

Se realizaron capturas directas de las microesferas empleando el sistema CMOS OV5647 en configuración de microscopía sin lentes.

Las imágenes adquiridas muestran estructuras circulares de bajo contraste, cercanas al límite de resolución impuesto por el tamaño de píxel ($\sim 1.4\,\mu m$).

\begin{figure}[H]
    \centering
    \includegraphics[width=0.5\textwidth]{fig/Fig_05_02_CMOS_RAW.png}
    \caption{Imagen RAW adquirida con el sensor CMOS en configuración lensless.}
    \label{fig:cmos_raw}
\end{figure}

Se observa que las microesferas no presentan bordes perfectamente definidos, sino patrones suavizados influenciados por fenómenos de difracción y por la respuesta espacial del sensor.
\section{Mejora perceptual mediante procesamiento computacional}

Con el objetivo de evaluar el impacto del procesamiento digital sobre la interpretabilidad de la imagen, se aplicaron técnicas de reconstrucción y super-resolución basadas en aprendizaje profundo.

\subsection{Reconstrucción FPM}

Se realizó una reconstrucción preliminar utilizando múltiples capturas bajo variación angular de iluminación.

\begin{figure}[H]
    \centering
    \includegraphics[width=0.5\textwidth]{fig/Fig_05_03_FPM.png}
    \caption{Resultado de reconstrucción mediante FPM.}
    \label{fig:fpm_result}
\end{figure}

Se observa mejora en contraste y definición visual, aunque la resolución física continúa limitada por el tamaño de píxel.

\subsection{Escalamiento mediante RealESRGAN}

Se aplicó super-resolución digital sobre la imagen base.

\begin{figure}[H]
    \centering
    \includegraphics[width=0.85\textwidth]{fig/Fig_05_04_SuperResolucion.png}
    \caption{Comparación entre imagen original y resultado tras super-resolución digital.}
    \label{fig:sr_result}
\end{figure}

Si bien la percepción visual mejora, se reconoce que estas técnicas no incrementan la resolución física real del sistema.
\section{Estimación experimental de la escala espacial}

Se procedió a realizar mediciones digitales del diámetro aparente de microesferas seleccionadas en regiones de interés (ROI).

El procedimiento consistió en:

\begin{itemize}
    \item Selección de microesferas aisladas.
    \item Medición del diámetro en píxeles.
    \item Conversión a unidades físicas mediante una escala preliminar ($\mu m/pixel$).
\end{itemize}

\begin{figure}[H]
    \centering
    \includegraphics[width=0.55\textwidth]{fig/Fig_05_05_Medicion_Diametro_CMOS.png}
    \caption{Ejemplo de medición digital del diámetro aparente en píxeles.}
    \label{fig:diametro_cmos}
\end{figure}

A partir de la comparación con el diámetro promedio obtenido por SEM, se estimó una escala espacial aproximada:

\[
\text{Escala experimental} \approx 0.17 \,\mu m/pixel
\]

La diferencia relativa entre mediciones se mantuvo dentro de un margen coherente considerando las limitaciones físicas del sistema.
\section{Análisis de limitaciones físicas y computacionales}

Los resultados experimentales evidencian que el sistema es capaz de detectar estructuras del orden de $2\,\mu m$, aunque operando cerca del límite de muestreo impuesto por el tamaño de píxel.

Las principales limitaciones identificadas incluyen:

\begin{itemize}
    \item Influencia significativa de la difracción en la definición de bordes.
    \item Dependencia crítica de la distancia muestra–sensor.
    \item Sensibilidad a variaciones en iluminación.
    \item Resolución efectiva limitada por el tamaño de píxel.
\end{itemize}

Estos factores explican las discrepancias observadas entre la imagen CMOS y las micrografías SEM, y refuerzan la necesidad de complementar la validación con patrones litográficos de mayor robustez geométrica.
\section{Síntesis de resultados y cumplimiento de objetivos}

Los resultados obtenidos permiten afirmar que:

\begin{itemize}
    \item Se implementó exitosamente un sistema de adquisición basado en sensor CMOS.
    \item Se estableció una metodología de validación metrológica preliminar.
    \item Se contrastaron mediciones con técnicas independientes (óptica y SEM).
    \item Se identificaron límites físicos asociados a resolución y difracción.
\end{itemize}

En consecuencia, el sistema demuestra coherencia dimensional preliminar y establece una base técnica sólida para futuras etapas de optimización del payload.

\chapter{Discusión}

\section{Coherencia dimensional del sistema CMOS}

Los resultados obtenidos evidencian que el sistema de adquisición basado en sensor CMOS es capaz de detectar estructuras del orden de $2\,\mu m$, aunque operando cerca de su límite físico de resolución.

La comparación con mediciones independientes realizadas mediante microscopía óptica calibrada y microscopía electrónica de barrido (SEM) permitió establecer una referencia dimensional externa. La diferencia relativa observada entre las mediciones digitales del sistema CMOS y la referencia SEM se mantiene dentro de un margen coherente con las limitaciones impuestas por el tamaño de píxel y los efectos difractivos.

Es importante destacar que la coherencia dimensional observada no implica que el sistema alcance resolución óptica equivalente a la de un microscopio convencional, sino que demuestra consistencia metrológica preliminar dentro de su régimen operativo.
\section{Influencia del tamaño de píxel y el límite de muestreo}

El sensor OV5647 posee un tamaño de píxel aproximado de $1.4\,\mu m$. Al intentar observar estructuras de $2\,\mu m$, el sistema opera en un régimen donde el objeto ocupa apenas unos pocos píxeles efectivos.

Desde el punto de vista del muestreo espacial, esto implica que la representación digital de la microesfera se encuentra cercana al límite de Nyquist, reduciendo la capacidad de definir bordes con nitidez.

La apariencia suavizada de las microesferas en las imágenes CMOS no constituye un error experimental, sino una consecuencia directa del muestreo discreto y de la respuesta espacial del detector.

Este comportamiento confirma que la resolución efectiva del sistema está dominada por:

\begin{itemize}
    \item Tamaño físico del píxel.
    \item Distancia muestra–sensor.
    \item Longitud de onda de la iluminación.
    \item Respuesta espectral del sensor.
\end{itemize}
\section{Efectos de difracción en configuración lensless}

En microscopía sin lentes, la formación de imagen no depende de un sistema óptico que enfoque la muestra, sino de la interacción directa entre la luz transmitida o dispersada y el plano del sensor.

Cuando las dimensiones del objeto son comparables a la longitud de onda de la iluminación o al tamaño de píxel, la imagen registrada se encuentra influenciada por fenómenos de difracción.

En el caso analizado, los patrones observados en las imágenes CMOS presentan características compatibles con un régimen dominado parcialmente por difracción, lo cual explica la transición suave entre objeto y fondo y la ausencia de bordes abruptos.

Este comportamiento es consistente con los fundamentos físicos de la formación de imagen en sistemas de contacto.
\section{Impacto del procesamiento computacional}

Las técnicas de reconstrucción y super-resolución aplicadas durante el análisis experimental mejoran la interpretabilidad visual de las estructuras observadas. Sin embargo, es fundamental diferenciar entre:

\begin{itemize}
    \item Mejora perceptual (contraste, nitidez aparente).
    \item Incremento real de resolución física.
\end{itemize}

Los algoritmos de reconstrucción no generan información física adicional más allá de la contenida en la señal original capturada por el sensor.

En consecuencia, si bien el procesamiento computacional constituye una herramienta valiosa para interpretación y análisis exploratorio, la validación metrológica debe fundamentarse en la imagen base adquirida bajo condiciones controladas.
\section{Limitaciones identificadas}

A partir de los resultados experimentales y su análisis teórico, se identifican las siguientes limitaciones principales del sistema:

\begin{itemize}
    \item Resolución efectiva limitada por tamaño de píxel.
    \item Sensibilidad crítica a la distancia muestra–sensor.
    \item Dependencia significativa de la estabilidad de iluminación.
    \item Influencia de ruido electrónico en regímenes de baja señal.
\end{itemize}

Estas limitaciones no invalidan la coherencia dimensional observada, pero establecen el marco físico dentro del cual deben interpretarse las mediciones.
\section{Implicancias para el desarrollo del payload}

Desde el punto de vista del proyecto INTISAT, los resultados obtenidos constituyen una validación preliminar de la viabilidad de integrar un sistema de microscopía computacional compacto como carga útil experimental.

La arquitectura lensless demuestra ventajas significativas en términos de:

\begin{itemize}
    \item Simplicidad mecánica.
    \item Bajo volumen y masa.
    \item Reducción de complejidad óptica.
\end{itemize}

No obstante, para aplicaciones de caracterización biológica más exigentes, será necesario considerar:

\begin{itemize}
    \item Uso de sensores con menor tamaño de píxel.
    \item Integración de patrones litográficos certificados para calibración robusta.
    \item Optimización del sistema de iluminación.
    \item Implementación controlada de técnicas avanzadas de reconstrucción.
\end{itemize}

La presente etapa debe entenderse como una validación fundacional que establece las bases técnicas para futuras iteraciones del sistema.

\chapter{Conclusiones y Trabajo Futuro}

\section{Conclusiones}

El presente trabajo tuvo como objetivo principal diseñar, implementar y validar preliminarmente un sistema de microscopía computacional basado en sensor CMOS, concebido como carga útil experimental del proyecto INTISAT.

A partir de los resultados obtenidos, se concluye que:

\begin{itemize}
    \item Se implementó exitosamente un sistema de adquisición de imágenes en configuración de microscopía sin lentes, empleando un sensor CMOS comercial y una arquitectura de control programático basada en Raspberry Pi 5.

    \item Se estableció una metodología de validación metrológica preliminar mediante comparación cruzada con microscopía óptica calibrada y microscopía electrónica de barrido (SEM), asegurando una referencia dimensional independiente.

    \item El sistema demostró capacidad para detectar estructuras del orden de $2\,\mu m$, operando cerca del límite impuesto por el tamaño de píxel del sensor.

    \item Las mediciones digitales obtenidas presentan coherencia dimensional preliminar al contrastarse con referencias externas, dentro de los márgenes esperables considerando efectos de muestreo y difracción.

    \item Se identificaron explícitamente las limitaciones físicas y computacionales del sistema, incluyendo restricciones de resolución efectiva, sensibilidad a la iluminación y dependencia crítica de la distancia muestra–sensor.

    \item El procesamiento computacional mejora la interpretabilidad visual, pero no incrementa la resolución física real del sistema, lo cual fue reconocido y documentado de manera conservadora.
\end{itemize}

En conjunto, el trabajo establece una base técnica sólida para el desarrollo progresivo del payload de microscopía computacional dentro del proyecto INTISAT.
\section{Trabajo Futuro}

Los resultados obtenidos abren múltiples líneas de desarrollo orientadas a mejorar la robustez metrológica y ampliar las capacidades del sistema.

Entre las principales líneas de trabajo futuro se consideran:

\begin{itemize}
    \item Incorporación de patrones litográficos certificados (grids micrométricos) para una calibración espacial más robusta y menos sensible a efectos difractivos.

    \item Evaluación de objetos patrón de mayor tamaño (5–10 µm) que permitan validar la escala espacial en un régimen menos cercano al límite de resolución.

    \item Optimización del sistema de iluminación, incluyendo caracterización espectral y análisis cuantitativo de uniformidad mediante herramientas fotométricas.

    \item Implementación controlada de técnicas avanzadas de reconstrucción computacional (por ejemplo, FPM) bajo un marco de validación física explícita.

    \item Evaluación de sensores CMOS con menor tamaño de píxel para incrementar la resolución efectiva sin comprometer la arquitectura compacta del sistema.

    \item Desarrollo de protocolos preliminares para observación exploratoria de muestras biológicas de tamaño micrométrico (por ejemplo, estructuras celulares ampliamente utilizadas en investigación básica), sin fines diagnósticos ni clínicos.
\end{itemize}

Desde la perspectiva del proyecto INTISAT, el sistema desarrollado representa una plataforma experimental modular que puede evolucionar progresivamente hacia aplicaciones más exigentes, manteniendo coherencia metrológica y trazabilidad experimental.

El presente trabajo constituye una etapa fundacional en el desarrollo del payload, estableciendo criterios físicos, metodológicos y documentales que permitirán consolidar futuras iteraciones del sistema.

\end{document}
